% 叩き台：解答の引数は空欄。問題本体は大問ごとに1ページ。
\MajorQuestionLesson{三人称単数現在の肯定文}
\TypeQuestion{英文を書き写し、日本語に訳しなさい。}{generalverb_third_affirmative_copy_translate}
\SetRowSpace{4mm}
\EngCopyTranslateSplit{He plays tennis.}{}
\EngCopyTranslateSplit{She studies English.}{}
\EngCopyTranslateSplit{Ken goes to school.}{}
\EngCopyTranslateSplit{My brother watches TV.}{}
\EngCopyTranslateSplit{She has a pen.}{}
\EngCopyTranslateSplit{They play tennis.}{}

\LessonPageBreak
\TypeQuestion{英文を書き写したあと、主語に下線を引き、動詞を丸で囲みなさい。}{generalverb_third_affirmative_copy_structure}
\SetRowSpace{3mm}
\EngWriteQuestion{\EngSentence{My sister likes music.}}{}
\EngWriteQuestion{\EngSentence{We study English.}}{}
\EngWriteQuestion{\EngSentence{Ken and Yui play tennis.}}{}
\EngWriteQuestion{\EngSentence{He has a book.}}{}
\EngWriteQuestion{\EngSentence{My brothers read books.}}{}
\EngWriteQuestion{\EngSentence{Yui goes to school.}}{}
\EngWriteQuestion{\EngSentence{She is busy.}}{}
\EngWriteQuestion{\EngSentence{My brother studies English.}}{}
\EngWriteQuestion{\EngSentence{They don't like tennis.}}{}
\EngWriteQuestion{\EngSentence{Ken does his homework.}}{}

\LessonPageBreak
\TypeQuestion{日本語に合うように、英文を必要なところだけ直して書きなさい。\par{\small ※短縮形の有無は、もとの英文に合わせること。}}{generalverb_third_affirmative_error_correct}
\SetRowSpace{3mm}
\EngWriteQuestion{\EngSentence{He play tennis.}\quad〔彼はテニスをします。〕}{}
\EngWriteQuestion{\EngSentence{She studys English.}\quad〔彼女は英語を勉強します。〕}{}
\EngWriteQuestion{\EngSentence{Ken watchs TV.}\quad〔ケンはテレビを見ます。〕}{}
\EngWriteQuestion{\EngSentence{My sister have a pen.}\quad〔私の姉はペンを持っています。〕}{}
\EngWriteQuestion{\EngSentence{They plays tennis.}\quad〔彼らはテニスをします。〕}{}
\EngWriteQuestion{\EngSentence{My brothers goes to school.}\quad〔私の兄たちは学校に行きます。〕}{}
\EngWriteQuestion{\EngSentence{Ken and Yui likes music.}\quad〔ケンとユイは音楽が好きです。〕}{}
\EngWriteQuestion{\EngSentence{She are busy.}\quad〔彼女は忙しいです。〕}{}
\EngWriteQuestion{\EngSentence{My brother do his homework.}\quad〔私の兄は宿題をします。〕}{}
\EngWriteQuestion{\EngSentence{I has a book.}\quad〔私は本を持っています。〕}{}

\LessonPageBreak
\TypeQuestion{日本語に合うように、（　）内の適切な語句を丸で囲み、完成した正しい英文を書きなさい。}{generalverb_third_affirmative_fill_choice}
\SetRowSpace{2mm}
\EngWriteQuestion{彼女は毎日英語を勉強します。\par\EngSentence{She \EngChoice{studys / studies / study} English every day.}}{}
\EngWriteQuestion{私たちはテニスをします。\par\EngSentence{We \EngChoice{plays / play} tennis.}}{}
\EngWriteQuestion{私の兄はテレビを見ます。\par\EngSentence{My brother \EngChoice{watch / watches} TV.}}{}
\EngWriteQuestion{ケンとユイは学校に行きます。\par\EngSentence{Ken and Yui \EngChoice{go / goes} to school.}}{}
\EngWriteQuestion{彼は本を持っています。\par\EngSentence{He \EngChoice{haves / have / has} a book.}}{}
\EngWriteQuestion{私の姉たちは音楽が好きです。\par\EngSentence{My sisters \EngChoice{likes / like} music.}}{}
\EngWriteQuestion{彼女は親切です。\par\EngSentence{She \EngChoice{is / does} kind.}}{}
\EngWriteQuestion{彼は毎日宿題をします。\par\EngSentence{He \EngChoice{do / does} his homework every day.}}{}

\LessonPageBreak
% レビュー反映：三単現を問う「すべて使用・1語不要・1語不足」とandを並べる問題を混在。
\TypeQuestion{日本語に合う英文を作りなさい。〔　〕内の語句は各問の指示に従って使うこと。\par{\small ※人名は日本語順、時の語句は文末に。}}{generalverb_third_affirmative_reorder}
\SetRowSpace{1.5mm}
\EngWriteQuestion{彼はテニスをします。\quad〔1語不要〕\par\EngReorder{tennis / plays / he / play}}{}
\EngWriteQuestion{私の姉は英語を勉強します。\quad〔すべて使用〕\par\EngReorder{English / studies / my sister}}{}
\EngWriteQuestion{ケンはテレビを見ます。\quad〔1語不足〕\par\EngReorder{TV / Ken}}{}
\EngWriteQuestion{彼女は本を持っています。\quad〔1語不足〕\par\EngReorder{a book / she}}{}
\EngWriteQuestion{私の兄たちは学校に行きます。\quad〔1語不要〕\par\EngReorder{goes / to school / my brothers / go}}{}
\EngWriteQuestion{ケンとユイは音楽が好きです。\quad〔すべて使用〕\par\EngReorder{Yui / like / Ken / music / and}}{}
\EngWriteQuestion{彼女は親切です。\quad〔すべて使用〕\par\EngReorder{kind / is / she}}{}
\EngWriteQuestion{彼は毎日宿題をします。\quad〔1語不足〕\par\EngReorder{his homework / he / every day}}{}

\LessonPageBreak
\TypeQuestion{英文の右の指示に従って、文を書き換えなさい。\par{\small ※主語に合わせ、動詞や所有格も必要なら直すこと。}}{generalverb_third_affirmative_transform}
\SetRowSpace{3mm}
\EngWriteQuestion{\EngSentence{I play tennis.}\quad〔主語を\EngText{He}に〕}{}
\EngWriteQuestion{\EngSentence{We study English.}\quad〔主語を\EngText{My sister}に〕}{}
\EngWriteQuestion{\EngSentence{She watches TV.}\quad〔主語を\EngText{They}に〕}{}
\EngWriteQuestion{\EngSentence{I have a book.}\quad〔主語を\EngText{Ken}に〕}{}
\EngWriteQuestion{\EngSentence{He goes to school.}\quad〔主語を\EngText{We}に〕}{}
\EngWriteQuestion{\EngSentence{They like music.}\quad〔主語を\EngText{Yui}に〕}{}
\EngWriteQuestion{\EngSentence{My brother reads books.}\quad〔主語を\EngText{My brothers}に〕}{}
\EngWriteQuestion{\EngSentence{Ken plays tennis.}\quad〔主語を\EngText{Ken and Yui}に〕}{}
\EngWriteQuestion{\EngSentence{They are busy.}\quad〔主語を\EngText{She}に〕}{}
\EngWriteQuestion{\EngSentence{You do your homework.}\quad〔主語を\EngText{He}に〕}{}

\WordBank{三人称単数現在の肯定文}{
  \WB{I / you}{代名詞}{私は／あなた（たち）は}{\EngText{I}は必ず大文字}
  \WB{we / they}{代名詞}{私たちは／彼らは}{\EngText{they}は物の複数にも使う}
  \WB{he / she / it}{代名詞}{彼は／彼女は／それは}{}
  \WB{my / your}{所有格}{私の／あなた（たち）の}{}
  \WB{his}{所有格}{彼の}{}
  \WB{a}{冠詞}{1つの／1人の}{訳さない場合もある}
  \WB{am / are / is}{be動詞}{～です／～にいます}{主語に合わせて使い分ける}
  \WB{do}{動詞等}{する；疑問・否定を作る}{三単現 \EngText{does}}
  \WB{does}{動詞等}{する；疑問・否定を作る}{\EngText{do}の三単現}
  \WB{don't}{短縮形}{しない；いいえ、しない}{\EngText{do not}の短縮形}
  \WB{play}{動詞}{（スポーツを）する}{三単現 \EngText{plays}}
  \WB{study}{動詞}{勉強する}{三単現 \EngText{studies}}
  \WB{like}{動詞}{～が好きである}{三単現 \EngText{likes}}
  \WB{read}{動詞}{読む}{三単現 \EngText{reads}}
  \WB{have}{動詞}{持っている；食べる}{三単現 \EngText{has}}
  \WB{go}{動詞}{行く}{三単現 \EngText{goes}}
  \WB{watch}{動詞}{見る}{三単現 \EngText{watches}}
  \WB{Ken / Yui}{人名}{ケン／ユイ}{最初の文字は大文字}
  \WB{brother}{名詞}{兄；弟}{複数形 \EngText{brothers}}
  \WB{sister}{名詞}{姉；妹}{複数形 \EngText{sisters}}
  \WB{English}{名詞}{英語}{最初の文字は大文字}
  \WB{music}{名詞}{音楽}{}
  \WB{tennis}{名詞}{テニス}{}
  \WB{book}{名詞}{本}{複数形 \EngText{books}}
  \WB{pen}{名詞}{ペン}{}
  \WB{school}{名詞}{学校}{}
  \WB{homework}{名詞}{宿題}{}
  \WB{TV}{名詞}{テレビ}{\EngText{television}の略}
  \WB{day}{名詞}{日；1日}{}
  \WB{busy}{形容詞}{忙しい}{}
  \WB{kind}{形容詞}{親切な}{}
  \WB{every}{形容詞}{毎～}{}
  \WB{to}{前置詞}{～へ／～に}{}
  \WB{and}{接続詞}{～と…}{}
  \WB{every day}{熟語}{毎日}{}
}

\PhraseBank{三人称単数現在の肯定文}{
  \PB{He plays tennis.}{彼はテニスをします。}
  \PB{She studies English.}{彼女は英語を勉強します。}
  \PB{Ken watches TV.}{ケンはテレビを見ます。}
  \PB{Yui goes to school.}{ユイは学校に行きます。}
  \PB{She has a pen.}{彼女はペンを持っています。}
  \PB{He does his homework.}{彼は宿題をします。}
}
